\documentclass[12pt]{article}

\usepackage[spanish]{babel}
\usepackage[utf8]{inputenc}
\usepackage{amsmath, amssymb}
\usepackage{geometry, float, fancyhdr, enumitem}

\geometry{letterpaper, margin=2.5cm}
\usepackage{fancyhdr}

\pagestyle{fancy}
\fancyhf{}

\fancyhead[L]{Diana Laura Diaz Juarez}
\fancyhead[R]{Boleta: 2024630047}

\renewcommand{\headrulewidth}{0.4pt}

\begin{document}

\begin{titlepage}
    \centering
    {\LARGE \scshape Instituto Politécnico Nacional} \par \vspace{1cm}
    {\Large \scshape Escuela Superior de Cómputo} \par \vspace{1.5cm}
    {\huge \bfseries 1er Lista de problemas y ejercicios} \par \vspace{2cm}
    {\Large \scshape Alumna: Diana Laura Diaz Juarez} \par \vspace{0.5cm}
    {\large \scshape Profesor: Gabriel Hurtado Áviles} \par \vspace{2cm}
    {\large Fecha: \today} \par
\begin{center}
    {\large{4CM4}\\}
\end{center}
\end{titlepage}

\section{Introducción}
Esta lista de problemas y ehercicios debe entregarse en formato PDF y preferentemente a computadora.
En la parte superior de cada hoja debe aparecer la informacion general del alumno (nombre completo, boeta y grupo).
Ademas, hay que incluir instrucchiones y/o enunciados de cada ejercicio, no solo las respuestas.
Finalmente, todas las respuestas de los ejecicios deben ir acompañadas del procedimiento y/o razonamiento.

\section{Problemas de estructuras ordenadas}
\subsection{Sea $\Sigma = \{0,1\}$. En los problemas 1-4, encuntre una cadena w que pertenezca al conjunto A con las caracteristicas indicadas, cuando sea posible.}


\textbf{A = $\Sigma^*$ Encuentre w tal que $|w| = 5$ }
El lenguaje $\Sigma^*$ contiene todas las cadenas formadas con los
símbolos $0$ y $1$, incluyendo $\lambda$.

Necesitamos una cadena de longitud 5:

\[
w = 01010
\]

Su longitud es:

\[
|w| = 5
\]

\textbf{Respuesta:}

\[
\boxed{w=01010}
\]

\textbf{A = $\Sigma^+$ Encuentre w tal que $|w| =0$}

\textbf{Procedimiento:}

La única cadena de longitud cero es la cadena vacía:

\[
\lambda
\]

Sin embargo:

\[
\Sigma^+ = \Sigma^*-\{\lambda\}
\]

Por lo tanto:

\[
\lambda \notin \Sigma^+
\]

\textbf{Respuesta:}

\[
\boxed{\text{No es posible}}
\]
\textbf{A = $\Sigma^*$ Encuentre w tal que $|w| =0$}
La cadena vacía $\lambda$ pertenece a $\Sigma^*$ y tiene longitud cero:

\[
|\lambda|=0
\]

\textbf{Respuesta:}

\[
\boxed{w=\lambda}
\]


\textbf{A = $\Sigma^5$ Encuentre w tal que el primer caracter de w sea 0}
Una cadena perteneciente a $\Sigma^5$ debe tener exactamente
5 símbolos.

Podemos elegir:

\[
w=01101
\]

Se cumple que:

\[
|w|=5
\]

y su primer carácter es $0$.

\textbf{Respuesta:}

\[
\boxed{w=01101}
\]


\subsection{Sea $\Sigma = \{a,b\}$, L1= {$ w \in \Sigma $ : w comienza con a} y L2 ={{ $w \in \Sigma^*$ : w termina en b}}.
Dadas las cadenas w=abbbaa, $\pi$ =aaabb y $\tau$=bbaba, realice los siguientes ejercicios 5-13:}

\textbf{Determine  $|\pi|$.}
La cadena es:

\[
\pi=aaabb
\]

Contamos sus símbolos:

\[
a\quad a\quad a\quad b\quad b
\]

Por lo tanto:

\[
|\pi|=5
\]

\textbf{Respuesta:}

\[
\boxed{|\pi|=5}
\]


\textbf{Encuentre $\omega\pi$.}
Tenemos:

\[
\omega=abbbaa
\]

y:

\[
\pi=aaabb
\]

La concatenación consiste en colocar $\pi$ después de $\omega$:

\[
\omega\pi=(abbbaa)(aaabb)
\]

Por lo tanto:

\[
\boxed{\omega\pi=abbbaaaaabb}
\]

\textbf{Encuentre $\pi\omega$.}
Colocamos $\omega$ después de $\pi$:

\[
\pi\omega=(aaabb)(abbbaa)
\]

Entonces:

\[
\boxed{\pi\omega=aaabbabbbaa}
\]

\textbf{Verifique si: \[ (\omega\pi)\tau=\omega(\pi\tau)\]}
Primero:

\[
\omega\pi=abbbaaaaabb
\]

Entonces:

\[
(\omega\pi)\tau
=(abbbaaaaabb)(baaba)
\]

\[
(\omega\pi)\tau=abbbaaaaabbbaaba
\]

Ahora:

\[
\pi\tau=(aaabb)(baaba)
\]

\[
\pi\tau=aaabbbaaba
\]

Por lo tanto:

\[
\omega(\pi\tau)
=(abbbaa)(aaabbbaaba)
\]

\[
\omega(\pi\tau)=abbbaaaaabbbaaba
\]

Como ambos resultados son iguales:

\[
\boxed{(\omega\pi)\tau=\omega(\pi\tau)}
\]

Por lo tanto, la igualdad es verdadera.



\textbf{Proporcione una palabra en $\Sigma^7$.}
La palabra debe estar formada por símbolos de $\Sigma=\{a,b\}$
y tener longitud 7.

Elegimos:

\[
w=abababa
\]

Como:

\[
|w|=7
\]

entonces:

\[
\boxed{w=abababa}
\]

\textbf{Proporcione una palabra en $L_1\cap\Sigma^5$.}
La palabra debe tener longitud 5 y comenzar con $a$.

Elegimos:

\[
w=ababa
\]

Se cumple:

\[
|w|=5
\]

y comienza con $a$.

Por lo tanto:

\[
\boxed{w=ababa}
\]


\textbf{Proporcione una palabra en $L_2\cap\Sigma^4$.}
La palabra debe tener longitud 4 y terminar en $b$.

Elegimos:

\[
w=aaab
\]

Se cumple:

\[
|w|=4
\]

y termina en $b$.

\textbf{Respuesta:}

\[
\boxed{w=aaab}
\]

\textbf{Proporcione una palabra en $L_1\cap L_2$.}
La palabra debe comenzar con $a$ y terminar con $b$.

Una opción es:

\[
w=ab
\]

Por lo tanto:

\[
\boxed{w=ab}
\]


\textbf{Determine $L_1\cap\Sigma^0$.}
Sabemos que:

\[
\Sigma^0=\{\lambda\}
\]

La cadena $\lambda$ no comienza con $a$, por lo que:

\[
\lambda\notin L_1
\]

Entonces:

\[
\boxed{L_1\cap\Sigma^0=\varnothing}
\]

\subsection{Sea $\Sigma= \{a,b,c,d\}$ y w=abbca, $\pi =abb$, $\tau =ca$. Consideremos los siguientes lenguajes:}

L1 = {$p \in \Sigma^*$ : $|p| \leq 5$}.

L2 = {$a^n b^m $: $n>m$}.

L3 = {$p \in \Sigma^*$ : p tiene mas letras b que letras a}.

Indique si las siguientes afirmaciones (ejercicios 14-46) son verdaderas o falsas1:



\textbf{$\lvert\omega\rvert=3$}
\[
\omega=abbca
\]
La cadena tiene cinco símbolos:
\[
\lvert\omega\rvert=5
\]
Por lo tanto:
\[
5\neq3
\]

\textbf{Respuesta:} Falso.


\textbf{$\lvert\omega\rvert=5$}
\[
\omega=abbca
\]
Por lo tanto:
\[
\lvert\omega\rvert=5
\]

\textbf{Respuesta:} Verdadero.

\textbf{$\omega=\tau\pi$ }
\[
\tau\pi=(ca)(abb)=caabb
\]
Mientras que:
\[
\omega=abbca
\]
Como:
\[
caabb\neq abbca
\]

\textbf{Respuesta:} Falso.

 \textbf{$\omega=\pi\tau$}
 \[
\pi\tau=(abb)(ca)=abbca
\]
Y:
\[
\omega=abbca
\]

\textbf{Respuesta:} Verdadero.

\textbf{$\omega(\pi\tau)=(\omega\pi)\tau$}
La concatenación de cadenas es asociativa:
\[
\omega(\pi\tau)=(\omega\pi)\tau
\]

\textbf{Respuesta:} Verdadero.

\textbf{$\omega(\tau\pi)=(\omega\pi)\tau$}
\[
\tau\pi=(ca)(abb)=caabb
\]
Entonces:
\[
\omega(\tau\pi)=abbcacaabb
\]

Por otro lado:
\[
(\omega\pi)\tau=(abbca)(abb)(ca)
\]
\[
(\omega\pi)\tau=abbcaabbca
\]

Como:
\[
abbcacaabb\neq abbcaabbca
\]

\textbf{Respuesta:} Falso.

\textbf{$\lambda\in\Sigma^+$}
La clausura positiva $\Sigma^+$ contiene cadenas de longitud mayor o igual a uno.
La cadena vacía tiene longitud cero:
\[
|\lambda|=0
\]
Por lo tanto:
\[
\lambda\notin\Sigma^+
\]

\textbf{Respuesta:} Falso.

\textbf{$\lambda\in\Sigma^*$}
La clausura de Kleene contiene la cadena vacía:
\[
\lambda\in\Sigma^*
\]

\textbf{Respuesta:} Verdadero.

\textbf{$\omega\in\Sigma^5$ }
Tenemos:
\[
\omega=abbca
\]
y:
\[
|\omega|=5
\]

Además, todos sus símbolos pertenecen a $\Sigma$.

\textbf{Respuesta:} Verdadero.

\textbf{$\omega\in\Sigma^3$}
\[
|\omega|=5
\]
pero $\Sigma^3$ contiene cadenas de longitud 3.

Por lo tanto:
\[
\omega\notin\Sigma^3
\]

\textbf{Respuesta:} Falso.

\textbf{$\pi\tau\in\Sigma^5$}
\[
|\pi|=3,\qquad|\tau|=2
\]

Entonces:
\[
|\pi\tau|=|\pi|+|\tau|=3+2=5
\]

\textbf{Respuesta:} Verdadero.

\textbf{$\pi\lambda\in\Sigma^3$}
La cadena vacía es el elemento neutro de la concatenación:
\[
\pi\lambda=\pi
\]

Como:
\[
|\pi|=3
\]

entonces:
\[
\pi\lambda\in\Sigma^3
\]

\textbf{Respuesta:} Verdadero.

\textbf{$a^8\omega\in\Sigma^8$}
\[
|a^8|=8
\]
y:
\[
|\omega|=5
\]

Por lo tanto:
\[
|a^8\omega|=8+5=13
\]

Como:
\[
13\neq8
\]

\textbf{Respuesta:} Falso.

\textbf{$a^3\omega\in\Sigma^8$}
\[
|a^3|=3
\]
y:
\[
|\omega|=5
\]

Entonces:
\[
|a^3\omega|=3+5=8
\]

\textbf{Respuesta:} Verdadero.

\textbf{$a^3c^3\in\Sigma^9$}
\[
|a^3c^3|=3+3=6
\]

Como:
\[
6\neq9
\]

\textbf{Respuesta:} Falso.

\textbf{$|a^2b^3|=|c^6|$}
\[
|a^2b^3|=2+3=5
\]

Mientras que:
\[
|c^6|=6
\]

Como:
\[
5\neq6
\]

\textbf{Respuesta:} Falso.

\textbf{$|a^4b^3c|=|a^3b^2c^3|$}
\[
|a^4b^3c|=4+3+1=8
\]

Y:
\[
|a^3b^2c^3|=3+2+3=8
\]

Por lo tanto:
\[
8=8
\]

\textbf{Respuesta:} Verdadero.

\textbf{$ab\in L_1$}
\[
L_1=\{\rho\in\Sigma^*:|\rho|\leq5\}
\]

Como:
\[
|ab|=2
\]

y:
\[
2\leq5
\]

\textbf{Respuesta:} Verdadero.

\textbf{$ab\in L_2$ }
\[
L_2=\{a^nb^m:n<m\}
\]

Para $ab$:
\[
n=1,\qquad m=1
\]

Pero:
\[
1\not<1
\]

\textbf{Respuesta:} Falso.

\textbf{$ab\in L_3$ }
\[
L_3=\{\rho\in\Sigma^*:\rho\text{ tiene más letras }b\text{ que letras }a\}
\]

En $ab$ tenemos:
\[
\#a=1,\qquad\#b=1
\]

No hay más letras $b$ que $a$.

\textbf{Respuesta:} Falso.

\textbf{$abab^2\in L_1$ }
\[
abab^2=ababb
\]

Por lo tanto:
\[
|abab^2|=5
\]

Como:
\[
5\leq5
\]

\textbf{Respuesta:} Verdadero.

\textbf{$abab^2\in L_2$ }
Una cadena de $L_2$ debe tener la forma:
\[
a^nb^m
\]

La cadena:
\[
ababb
\]
alterna entre $a$ y $b$, por lo que no tiene esa forma.

\textbf{Respuesta:} Falso.

\textbf{$abab^2\in L_3$}
\[
abab^2=ababb
\]

Tiene:
\[
\#a=2,\qquad\#b=3
\]

Como:
\[
3>2
\]

\textbf{Respuesta:} Verdadero.

\textbf{$abcad\in L_1$}
\[
|abcad|=5
\]

Como:
\[
5\leq5
\]

\textbf{Respuesta:} Verdadero.

\textbf{$abcad\in L_2$}
Las cadenas de $L_2$ tienen la forma:
\[
a^nb^m
\]

La cadena $abcad$ contiene los símbolos $c$ y $d$, por lo que no tiene esa forma.

\textbf{Respuesta:} Falso.

\textbf{$abcad\in L_3$}
En:
\[
abcad
\]

tenemos:
\[
\#a=2,\qquad\#b=1
\]

No hay más $b$ que $a$.

\textbf{Respuesta:} Falso.

\textbf{$L_1\subseteq L_2$}
Tomemos:
\[
a\in L_1
\]

porque:
\[
|a|=1\leq5
\]

Pero:
\[
a\notin L_2
\]

porque una cadena de $L_2$ debe contener $a$ y $b$ cumpliendo $n<m$.

Por lo tanto:
\[
L_1\nsubseteq L_2
\]

\textbf{Respuesta:} Falso.

\textbf{$L_1\subseteq L_3$ }
Tomemos:
\[
a\in L_1
\]

Pero $a$ no tiene más letras $b$ que letras $a$.

Por lo tanto:
\[
a\notin L_3
\]

Entonces:
\[
L_1\nsubseteq L_3
\]

\textbf{Respuesta:} Falso.

\textbf{$L_2\subseteq L_3$}
Si:
\[
w\in L_2
\]

entonces:
\[
w=a^nb^m
\]

con:
\[
n<m
\]

Por lo tanto, $w$ tiene más letras $b$ que letras $a$.

Entonces:
\[
w\in L_3
\]

y por lo tanto:
\[
L_2\subseteq L_3
\]

\textbf{Respuesta:} Verdadero.

\textbf{$L_3\subseteq L_2$}
Consideremos:
\[
ba
\]

Tiene:
\[
\#a=1,\qquad\#b=1
\]

Por lo tanto, en realidad $ba\notin L_3$.

Tomemos mejor:
\[
bba
\]

Entonces:
\[
\#b=2,\qquad\#a=1
\]

por lo que:
\[
bba\in L_3
\]

Sin embargo, $bba$ no tiene la forma:
\[
a^nb^m
\]

Por lo tanto:
\[
bba\notin L_2
\]

\[
L_3\nsubseteq L_2
\]

\textbf{Respuesta:} Falso.

\textbf{$\alpha\in\Sigma^*$, entonces $alpha\lambda=\lambda\alpha$}

La cadena vacía $\lambda$ es el elemento neutro de la concatenación:
\[
\alpha\lambda=\alpha
\]

y:
\[
\lambda\alpha=\alpha
\]

Por lo tanto:
\[
\alpha\lambda=\lambda\alpha
\]

\textbf{Respuesta:} Verdadero.

\textbf{Si $\alpha\in\Sigma^3$ y $\beta\in\Sigma^4$, entonces $|\alpha\beta|=12$}
\[
|\alpha|=3,\qquad|\beta|=4
\]

Entonces:
\[
|\alpha\beta|=|\alpha|+|\beta|
\]

\[
|\alpha\beta|=3+4=7
\]

Como:
\[
7\neq12
\]

\textbf{Respuesta:} Falso.

\textbf{Si $\alpha\in\Sigma^3$ y $\beta\in\Sigma^4$, entonces Si $\alpha\in\Sigma^3$ y $|\alpha\beta|=7$}
\[
|\alpha\beta|=|\alpha|+|\beta|
\]

\[
|\alpha\beta|=3+4=7
\]

\textbf{Respuesta:} Verdadero.

\subsection{En caso de que sea posible, encuentre cinco palabras de cada uno de los lenguajes que se definen en los problemas 47-56, usando el alfabeto $\Sigma = \{0,1,2,3,4,5,6\}$}.

\textbf{$L_1=\{01^n:n\in\mathbb{N}\}$.}
La expresión $01^n$ indica que toda palabra comienza con $0$ y después
contiene $n$ unos.

Por ejemplo:
\[
0,\quad 01,\quad 011,\quad 0111,\quad \ldots
\]

\textbf{Respuesta:}
\[
L_1=\{0,01,011,0111,\ldots\}
\]

\textbf{$L_2=\{0^n1^n2^n:n\in\mathbb{N}\}$.}
Las palabras deben tener la misma cantidad de $0$, $1$ y $2$.

Para $n=1$:
\[
012
\]

Para $n=2$:
\[
001122
\]

Para $n=3$:
\[
000111222
\]

\textbf{Respuesta:}
\[
L_2=\{012,001122,000111222,\ldots\}
\]

\textbf{$L_3=\{\omega\in\Sigma^*:|\omega|=6\}.$}
Son todas las palabras sobre $\Sigma$ que tienen exactamente seis símbolos.

\textbf{Respuesta:}
\[
L_3=\{\omega\in\Sigma^*:|\omega|=6\}.
\]

Además,
\[
|L_3|=|\Sigma|^6=7^6=117649.
\]

\textbf{$L_4=\{\omega\in\Sigma^*:\text{la suma de los dígitos de }\omega
\text{ es }6\}.$}
Se buscan palabras cuya suma de símbolos sea exactamente $6$.

Algunos ejemplos son:
\[
000006,\quad 000015,\quad 000024,\quad
000033,\quad 000042.
\]

\textbf{Respuesta:}

\[
L_4=\{\omega\in\Sigma^*:\operatorname{suma}(\omega)=6\}.
\]

\textbf{Determine $L_1\cap L_2$.}
Una palabra debe pertenecer simultáneamente a
\[
L_1=\{01^n:n\in\mathbb N\}
\]
y
\[
L_2=\{0^n1^n2^n:n\in\mathbb N\}.
\]

En $L_1$ solamente aparecen $0$ y $1$, mientras que en $L_2$
aparece también al menos un $2$.

Por lo tanto, no existe una palabra que pertenezca a ambos lenguajes.

\textbf{Respuesta:}
\[
\boxed{L_1\cap L_2=\varnothing}
\]

\textbf{Determine $L_1\cap L_3$.}
Las palabras de $L_3$ tienen longitud $6$.

En $L_1$ las palabras tienen la forma $01^n$.

Para que tengan longitud $6$:
\[
1+n=6
\]
\[
n=5.
\]

Por lo tanto, la única palabra es:
\[
011111.
\]

\textbf{Respuesta:}
\[
\boxed{L_1\cap L_3=\{011111\}}
\]

\textbf{Determine $L_2\cap L_3$.}
En $L_3$ las palabras tienen longitud $6$.

Para una palabra de $L_2$:
\[
|0^n1^n2^n|=3n.
\]

Entonces:
\[
3n=6
\]
\[
n=2.
\]

La palabra correspondiente es:
\[
001122.
\]

\textbf{Respuesta:}
\[
\boxed{L_2\cap L_3=\{001122\}}
\]

\textbf{Determine $L_1\cap L_4$.}
Las palabras de $L_1$ tienen la forma:
\[
01^n.
\]

La suma de sus símbolos es:
\[
0+n=n.
\]

Para pertenecer a $L_4$:
\[
n=6.
\]

Entonces:
\[
01^6=0111111.
\]

\textbf{Respuesta:}
\[
\boxed{L_1\cap L_4=\{0111111\}}
\]

\textbf{Determine $L_2\cap L_4$.}
Una palabra de $L_2$ tiene la forma:
\[
0^n1^n2^n.
\]

La suma de sus símbolos es:
\[
0n+1n+2n=3n.
\]

Para pertenecer a $L_4$:
\[
3n=6.
\]

Por lo tanto:
\[
n=2.
\]

La palabra es:
\[
001122.
\]

\textbf{Respuesta:}
\[
\boxed{L_2\cap L_4=\{001122\}}
\]

\textbf{Determine $L_3\cap L_4$.}
La palabra debe tener longitud $6$ y la suma de sus dígitos debe ser $6$.

Algunos elementos son:
\[
000006,\quad
000015,\quad
000024,\quad
000033,\quad
000042.
\]

\textbf{Respuesta:}

\[
\boxed{
L_3\cap L_4=
\{\omega\in\Sigma^*:|\omega|=6
\text{ y la suma de sus dígitos es }6\}
}
\]

\section{Propiedades de la clausura de kleene y positiva}

\section{Reflexion o inverso de un lenguaje}

\subsection{Dado A un lenguaje sobre $\Sigma$ se define $A^R$ de la siguiente forma:}
 
\textbf{$ (A\cdot B)^R=B^R\cdot A^R.$}
Sea $w\in(A\cdot B)^R$.

Entonces existe $u\in A$ y $v\in B$ tales que:
\[
w=(uv)^R.
\]

Por la definición de reversa:
\[
(uv)^R=v^Ru^R.
\]

Como
\[
v^R\in B^R
\]
y
\[
u^R\in A^R,
\]
entonces:
\[
w\in B^RA^R.
\]

Por lo tanto:
\[
\boxed{(A\cdot B)^R=B^R\cdot A^R}.
\]


\textbf{($A\cup B)^R=A^R\cup B^R.$}
Si $w\in A\cup B$, entonces:
\[
w\in A \quad\text{o}\quad w\in B.
\]

Por lo tanto:
\[
w^R\in A^R \quad\text{o}\quad w^R\in B^R.
\]

Así:
\[
w^R\in A^R\cup B^R.
\]

\textbf{Respuesta:}
\[
\boxed{(A\cup B)^R=A^R\cup B^R}
\]

\textbf{($A\cap B)^R=A^R\cap B^R.$}
Si
\[
w\in A\cap B,
\]
entonces:
\[
w\in A
\quad\text{y}\quad
w\in B.
\]

Al invertir:
\[
w^R\in A^R
\quad\text{y}\quad
w^R\in B^R.
\]

Por lo tanto:
\[
w^R\in A^R\cap B^R.
\]

\textbf{Respuesta:}
\[
\boxed{(A\cap B)^R=A^R\cap B^R}
\]


\textbf{($A^R)^R=A.$}
Invertir una palabra dos veces devuelve la palabra original:
\[
(w^R)^R=w.
\]

Por lo tanto, al invertir dos veces todos los elementos de $A$,
obtenemos nuevamente $A$.

\textbf{Respuesta:}
\[
\boxed{(A^R)^R=A}
\]

\textbf{($A^*)^R=(A^R)^*.$}
Por definición:
\[
A^*=\bigcup_{n\geq0}A^n.
\]

Aplicando reversa:
\[
(A^*)^R=
\left(\bigcup_{n\geq0}A^n\right)^R.
\]

La reversa de una concatenación de palabras de $A$ produce
concatenaciones de palabras de $A^R$.

Por lo tanto:
\[
(A^*)^R=\bigcup_{n\geq0}(A^R)^n.
\]

Así:
\[
\boxed{(A^*)^R=(A^R)^*}.
\]

\textbf{($A^+)^R=(A^R)^+.$}
Por definición:
\[
A^+=\bigcup_{n\geq1}A^n.
\]

Aplicando reversa:
\[
(A^+)^R=
\bigcup_{n\geq1}(A^n)^R.
\]

Cada potencia invertida corresponde a una potencia de $A^R$.

Entonces:
\[
(A^+)^R=
\bigcup_{n\geq1}(A^R)^n.
\]

\textbf{Respuesta:}
\[
\boxed{(A^+)^R=(A^R)^+}
\]

\section{Longitud de cadenas}

\textbf{Si $w\in\Sigma^*$ y $n,m\in\mathbb N$, demostrar que
\[
|w^{n+m}|=|w^n|+|w^m|.
\]}

\[
w^{n+m}=w^nw^m.
\]

Usando la propiedad de la longitud de una concatenación:
\[
|uv|=|u|+|v|,
\]
se obtiene:
\[
|w^{n+m}|=|w^n|+|w^m|.
\]

\textbf{Respuesta:}
\[
\boxed{|w^{n+m}|=|w^n|+|w^m|}
\]

\section{n-cadenas o n-palabras}

\section{Prefijos y sufijos de una cadena}

\textbf{Enumere todos los prefijos de u}
Los prefijos se obtienen tomando los símbolos desde el inicio:

\[
\lambda
\]

\[
b
\]

\[
bc
\]

\[
bcb
\]

\[
bcba
\]

\[
bcbaa
\]

\[
bcbaad
\]

\[
bcbaadb.
\]

\textbf{Respuesta:}

\[
\boxed{
\operatorname{Pref}(u)=
\{\lambda,b,bc,bcb,bcba,bcbaa,bcbaad,bcbaadb\}
}
\]
\textbf{Enumere todos los sufijos de u}
Los sufijos se obtienen tomando los símbolos desde el final:

\[
\lambda
\]

\[
b
\]

\[
db
\]

\[
adb
\]

\[
aadb
\]

\[
baadb
\]

\[
cbaadb
\]

\[
bcbaadb.
\]

\textbf{Respuesta:}

\[
\boxed{
\operatorname{Suf}(u)=
\{\lambda,b,db,adb,aadb,baadb,cbaadb,bcbaadb\}
}
\]

\section{Potencias de un lenguaje}

\textbf{Sea
\[
\Sigma=\{0,1\}.
\]
Determinar $\Sigma^2$, $\Sigma^3$ y $\Sigma^4$.}
Para $\Sigma^2$:

\[
\Sigma^2=
\{00,01,10,11\}.
\]

Para $\Sigma^3$:

\[
\Sigma^3=
\{000,001,010,011,100,101,110,111\}.
\]

Para $\Sigma^4$:

\[
\Sigma^4=
\{
0000,0001,0010,0011,
0100,0101,0110,0111,
1000,1001,1010,1011,
1100,1101,1110,1111
\}.
\]

\textbf{Demostrar que
\[
|uv|=|u|+|v|.
\]}
La concatenación $uv$ coloca todos los símbolos de $u$ seguidos de
todos los símbolos de $v$.

Por lo tanto, el número total de símbolos es la suma de sus longitudes.

\textbf{Respuesta:}
\[
\boxed{|uv|=|u|+|v|}
\]

\textbf{Demostrar por inducción que
\[
|u^n|=n|u|.
\]}
Caso base: $n=0$.

\[
|u^0|=|\lambda|=0
\]
y
\[
0|u|=0.
\]

\textbf{Hipótesis:}

\[
|u^n|=n|u|.
\]

\textbf{Paso inductivo:}

\[
u^{n+1}=u^nu.
\]

Entonces:
\[
|u^{n+1}|
=
|u^nu|
=
|u^n|+|u|.
\]

Aplicando la hipótesis:
\[
|u^{n+1}|
=
n|u|+|u|
=
(n+1)|u|.
\]

\textbf{Respuesta:}
\[
\boxed{|u^n|=n|u|}
\]

\textbf{Demostrar que
\[
(uv)^R=v^Ru^R.
\]}
Por la definición recursiva de reversa, al invertir una palabra se cambia
el orden de sus símbolos.

Si primero tenemos $u$ y después $v$:
\[
uv,
\]
al invertirla, $v$ queda primero y $u$ después.

Por lo tanto:
\[
(uv)^R=v^Ru^R.
\]

\textbf{Respuesta:}
\[
\boxed{(uv)^R=v^Ru^R}
\]










\end{document}